\section{Experiments}
\label{sec:experiments}

\subsection{Setup}

We initialize from BART-large~\citep{lewis2020bart}, replace encoder self-attention with \model{}, and extend the position embeddings to 4{,}096 tokens by copying. Models are fine-tuned for 50{,}000 steps with a batch size of 8 and a peak learning rate of $3 \times 10^{-5}$. We report ROUGE~\citep{lin2004rouge} on the test sets of arXiv, PubMed, and GovReport.

\subsection{Main results}

Table~\ref{tab:main} shows that reading more of the document helps on every dataset, and that learned block selection is better than a fixed global pattern at the same budget.

\begin{table}[t]
  \centering
  \caption{ROUGE-L on the test sets. Memory is the peak during training at batch size 1.}
  \label{tab:main}
  \small
  \begin{tabular}{lcccc}
    \toprule
    Model & arXiv & PubMed & GovRep. & Mem. \\
    \midrule
    BART (1{,}024) & 39.8 & 40.1 & 50.2 & 11\,GB \\
    Full attn. (4{,}096) & 41.5 & 41.9 & 53.0 & 38\,GB \\
    Longformer & 41.2 & 41.6 & 52.4 & 24\,GB \\
    BigBird & 41.3 & 41.8 & 52.7 & 26\,GB \\
    \model{} (ours) & \textbf{41.6} & \textbf{42.3} & \textbf{53.4} & 23\,GB \\
    \bottomrule
  \end{tabular}
\end{table}

\subsection{Ablations}

Removing the learned selection and attending to $r$ random blocks costs 0.7 ROUGE-L on arXiv. Halving the local window to $w = 128$ costs a further 0.4 points, which suggests that local and global context are complementary.
